\chapter{FDE Support, Merge, and Glut-Free Reconstruction}
\label{app:fde-support-merge}

This appendix collects the finite support calculus used in Parts III and VI in one place. It adds no new philosophical claim; its purpose is to make explicit the definitions and elementary propositions on which Chapters 10, 12, 20, and 22 rely.

\section*{Signed support and the four values}
Fix a proposition $p$. Let positive and negative support be represented by two indicators
\[
\sigma^+(p),\sigma^-(p)\in\{0,1\}.
\]
Define the support-value map $q:\{0,1\}^2\to\mathbb F$ by
\[
q(0,0)=N,\qquad q(1,0)=T,\qquad q(0,1)=F,\qquad q(1,1)=B.
\]
Equip $\mathbb F=\{N,T,F,B\}$ with the information order
\[
N\sqsubseteq T\sqsubseteq B,\qquad N\sqsubseteq F\sqsubseteq B,
\]
with $T$ and $F$ incomparable. Thus $\mathbb F$ is a finite complete lattice and, in particular,
\[
T\vee F=B,\qquad N\vee x=x,\qquad B\vee x=B.
\]

The same construction may be made over evidential witness groups. If $\mathcal E^+(U)$ and $\mathcal E^-(U)$ are the positive and negative evidence groups on a region $U$, define
\[
q(e^+,e^-)=
\begin{cases}
N,&e^+=0,\ e^-=0,\\
T,&e^+\neq0,\ e^-=0,\\
F,&e^+=0,\ e^-\neq0,\\
B,&e^+\neq0,\ e^-\neq0.
\end{cases}
\]
The four-valued semantics is therefore read from two independent support coordinates rather than taken as primitive.

\section*{Information merge}
Let $s_1,\ldots,s_n$ be independent support reports about $p$. Define
\[
M(p)=\bigvee_{i=1}^n s_i(p).
\]
When the reports are indexed by patches of a cover, Chapter 22 writes the same operation as
\[
\Sigma(p)=\bigvee_i q_i(p).
\]
Hence $\Sigma=M$ after specialization to cover-indexed sources.

\paragraph{Proposition A.1 (Support--Merge).}
$M(p)=B$ if and only if the family collectively contains positive support and negative support for $p$. Moreover, $M(p)=N$ if and only if every report is $N$.

\paragraph{Proof.}
Write each report as a pair $(\sigma_i^+,\sigma_i^-)$. Join in the information order is coordinatewise Boolean disjunction, so
\[
M(p)=q\!\left(\bigvee_i\sigma_i^+,\ \bigvee_i\sigma_i^-\right).
\]
The result is $B$ exactly when both displayed coordinates are $1$, and is $N$ exactly when both are $0$. \hfill$\square$

\section*{Glut-free reconstruction}
Let $\mathbb F_{\neg B}=\{N,T,F\}$. A glut-free common extension of a family $\{s_i(p)\}$ is an element $u\in\mathbb F_{\neg B}$ such that $s_i(p)\sqsubseteq u$ for every $i$.

\paragraph{Proposition A.2 (Glut-Free Reconstruction Obstruction).}
A glut-free common extension exists if and only if $M(p)\neq B$.

\paragraph{Proof.}
If $M(p)\neq B$, then $M(p)\in\mathbb F_{\neg B}$ and, by the defining property of a join, every $s_i(p)\sqsubseteq M(p)$. Conversely, if $M(p)=B$, the family contains both polarities. Any common upper bound must therefore dominate both $T$ and $F$, whose least upper bound is $B$. No element of $\mathbb F_{\neg B}$ can do so. \hfill$\square$

This proposition makes the gap/glut asymmetry exact. $N$ is not an obstruction: $N\vee x=x$. A gap records insufficient polar evidence. A glut records jointly retained positive and negative support, and becomes an obstruction only relative to a reconstruction regime that excludes $B$.

\chapter{Covers, Coherence Defects, and Mutual Non-Determination}
\label{app:coherence-defects}

This appendix isolates the cover-level apparatus of Chapters 20--23. Its main purpose is to prevent three operations from being conflated:
\[
\text{restriction}\neq\text{gluing}\neq\text{merge/repair}.
\]

\section*{Evidence presheaves and restriction}
Let $X$ be a domain with cover $\mathcal U=\{U_i\}$. For each open region $U$, let $\mathcal E^+(U)$ and $\mathcal E^-(U)$ be abelian groups generated by positive and negative evidential witnesses. For $V\subseteq U$, let
\[
\rho_{V,U}^{\pm}:\mathcal E^{\pm}(U)\to\mathcal E^{\pm}(V)
\]
be restriction maps satisfying identity and composition. The intended reading in Part VI permits information loss under restriction: evidence available on $U$ need not remain reconstructible from $V$ alone.

For a chosen family $e_i^\pm\in\mathcal E^\pm(U_i)$, define the overlap defect
\[
(\delta^\pm e)_{ij}
=
\rho_{U_{ij},U_j}^{\pm}(e_j^\pm)
-
\rho_{U_{ij},U_i}^{\pm}(e_i^\pm),
\qquad U_{ij}=U_i\cap U_j.
\]
The family is polarity-coherent when $\delta^\pm e=0$.

\paragraph{Remark.}
$\delta e$ is a \v{C}ech $1$-cochain and, for a chosen $0$-cochain $e$, is a coboundary by construction. Consequently its cohomology class cannot diagnose whether these particular witnesses agree: $[\delta e]=0$ automatically whenever the ordinary cochain complex is defined. The relevant test here is the cochain-level equation $\delta e=0$, not nonvanishing of $H^1$.

\section*{Ordinary gluing versus information merge}
Ordinary sheaf gluing begins with a matching family,
\[
\rho_{U_{ij},U_i}(e_i)=\rho_{U_{ij},U_j}(e_j),
\]
and produces a global section restricting exactly to each local section. By contrast, the FDE information merge
\[
\Sigma(p)=\bigvee_i q_i(p)
\]
can combine reports that do not match. If one patch reports $T$ and another $F$, the merge returns $B$; it does not produce a section whose restrictions are simultaneously the original $T$ and $F$ reports. Merge is therefore an information-preserving repair operation, not ordinary sheaf gluing.

\section*{Mutual non-determination}
Let the semantic coordinate be the merged value $\Sigma(p)\in\{N,T,F,B\}$ and the coherence coordinate be whether $(\delta^+(p),\delta^-(p))$ vanishes.

\paragraph{Proposition B.1 (Mutual Non-Determination).}
Neither the merged FDE value nor the coherence status determines the other.

\paragraph{Proof by realizability.}
A two-patch cover suffices. For $\Sigma=T$, choose either a shared positive witness on the overlap (coherent) or two distinct positive witnesses that restrict differently (incoherent). The same construction with negative witnesses realizes coherent and incoherent cases with $\Sigma=F$. For $\Sigma=B$, choose positive and negative evidence whose respective witnesses agree across the overlap to obtain a coherent glut, or choose a mismatch in at least one polarity to obtain an incoherent glut. Thus fixing $\Sigma$ does not fix coherence. Conversely, coherent and incoherent configurations occur at several distinct merged values, so coherence does not fix $\Sigma$. \hfill$\square$

The proposition is the formal core of the book's claim that contradiction and obstruction are different coordinates. A glut can be coherent, and a classical-looking verdict can be evidentially incoherent.

\section*{Action-relative coherence liability}
Chapter 23 extends the repair functional from
\[
E(a)+\lambda R(a)
\]
to
\[
J(a)=E(a)+\lambda R(a)+\mu C(a),
\qquad
C(a,p)=\rho_a\bigl(\delta^+(p),\delta^-(p)\bigr).
\]
The dependence on $a$ is essential: witness mismatch can be corroborative for an action insensitive to witness identity and disqualifying for an action that requires identity of the underlying witness. Hence coherence defect is structural, while coherence liability is action-relative.

\chapter{Resolution Profiles and Dispersion}
\label{app:dispersion}

Part IV separates degree-valued semantics from FDE's evidential polarity and then adds resolution as an explicit parameter. This appendix records the minimal mathematical structure needed for that move.

\section*{Resolution profiles}
Let $(\Lambda,\preceq)$ index resolutions or scales, and let
\[
v_\lambda(p)\in V
\]
be the value assigned to proposition $p$ at resolution $\lambda$, where $V$ carries a distance or other quantitative comparison $d$. The value $v_\lambda(p)$ and the full profile
\[
\lambda\longmapsto v_\lambda(p)
\]
are distinct objects. Many-valued semantics alone supplies the former; interpreting the latter as resolution behavior is an additional modeling choice of this book.

For a finite refinement or coarsening chain
\[
\lambda_0\preceq\lambda_1\preceq\cdots\preceq\lambda_n,
\]
define local increments
\[
\epsilon_i=d\bigl(v_{\lambda_i}(p),v_{\lambda_{i+1}}(p)\bigr)
\]
and accumulated variation
\[
D_{0:n}(p)=\sum_{i=0}^{n-1}\epsilon_i.
\]

\paragraph{Proposition C.1 (Local Smallness Does Not Imply Global Persistence).}
For every $\varepsilon>0$ and every target displacement $L>0$, there are sufficiently long chains for which each $\epsilon_i<\varepsilon$ while $D_{0:n}(p)\ge L$.

\paragraph{Proof.}
Choose $n>L/\varepsilon$ and a profile with constant increment $L/n$. Then $L/n<\varepsilon$ while the accumulated variation is $n(L/n)=L$. \hfill$\square$

The proposition captures the mechanism used in Chapter 15's sorites analysis. No adjacent step need contain a contradiction or a coherence failure. Dispersion is instead the accumulation of individually small changes into macroscopic displacement.

\section*{Persistence and topology}
A persistence functional $\mathcal P[s]$ may be chosen to measure how much of a verdict survives along a refinement/coarsening trajectory. The manuscript deliberately leaves its canonical form open; one possible family of models takes $\mathcal P$ to decrease with accumulated variation $D_{0:n}$.

This requires more structure than a bare topology. A topology can specify neighborhoods and continuity, but does not by itself assign magnitudes to successive displacements or provide an additive accumulation law. Consequently the qualitative topology proposed in Chapter 24 and the dispersion mechanism of Chapter 15 are distinct but related: the former can organize nearness under refinement, while the latter additionally requires a metric, uniformity, filtration, persistence functional, or comparable quantitative structure.

The resulting separation is
\[
\boxed{\text{compatibility}\neq\text{continuity}\neq\text{persistence}.}
\]
Compatibility asks whether fixed local data agree. Continuity asks which configurations are nearby. Persistence asks whether repeated locally small displacement remains globally small.
